\chapter{The Independence Boundary}
\label{app:independence-boundary}

The main thesis concerns the regular weak null \(I(P)=I(Q)>0\). This appendix
explains why constructing an independent reference distribution does not turn
the Expanded Welch method into a new first-order test of independence.

\section{Constructing the independent reference}

For a joint distribution $P=(p_{ij})$, define
\begin{equation}
  q_{ij}=p_{i+}p_{+j}.
  \label{eq:appendix-independent-q}
\end{equation}
Then $Q=(q_{ij})$ is a probability distribution because
\begin{equation}
  \sum_{i,j}q_{ij}
  =\left(\sum_i p_{i+}\right)
   \left(\sum_j p_{+j}\right)=1.
\end{equation}
It has the same margins as \(P\), since
\begin{equation}
  \sum_jq_{ij}=p_{i+},
  \qquad
  \sum_iq_{ij}=p_{+j}.
\end{equation}
The factorisation in Equation~\eqref{eq:appendix-independent-q} makes \(X\)
and \(Y\) independent under \(Q\), so \(I(Q)=0\).

The divergence from $P$ to this reference is
\begin{align}
  D_{\mathrm{KL}}(P\Vert Q)
  &=\sum_{i,j}p_{ij}\log\frac{p_{ij}}{q_{ij}}\\
  &=\sum_{i,j}p_{ij}\log\frac{p_{ij}}{p_{i+}p_{+j}}\\
  &=I(P).
\end{align}
Consequently, testing \(I(P)=I(Q)\) for this constructed reference is exactly
testing \(I(P)=0\), or independence.

\section{Why the first-order term vanishes}

At independence, $p_{ij}=p_{i+}p_{+j}$, and every population pointwise-MI
value is zero:
\begin{equation}
  \pmi_P(i,j)=\log 1=0.
\end{equation}
Equation~\eqref{eq:mi-derivative} therefore gives
\begin{equation}
  \left.\frac{\mathrm d I(P_\varepsilon)}
  {\mathrm d\varepsilon}\right|_0
  =\pmi_P(x,y)-I(P)=0.
\end{equation}
The first-order variance also vanishes:
\begin{equation}
  V(P)=\Var_P\{\pmi_P(X,Y)\}=0.
\end{equation}
The denominator used by Normal Wald and Expanded Welch is thus
degenerate under the population independence null. This does not mean that a
finite sample must have zero plug-in MI or zero plug-in variance. Sampling
noise usually makes the empirical table depart from exact independence. The
point is that these random departures do not enter MI through the usual
first-order term at the null. They enter through the second-order term derived
next, so changing only the degrees of freedom of a first-order statistic cannot
supply the missing source of variation.

\section{Second-order behaviour}

Fix a strictly positive independent table $P$, and let $h=(h_{ij})$ describe a
small redistribution of probability, with $\sum_{i,j}h_{ij}=0$. Define the path
\begin{equation}
  P_\varepsilon=P+\varepsilon h,
  \qquad
  f(\varepsilon)=I(P_\varepsilon),
  \label{eq:independence-path}
\end{equation}
for values of $\varepsilon$ small enough that all cells remain positive. Here
$\varepsilon=0$ is the expansion point and $f(0)=I(P)=0$. The second-order
Taylor expansion is
\begin{equation}
  f(\varepsilon)
  \approx
  \underbrace{f(0)}_{\text{zeroth order}}
  +\underbrace{f'(0)\varepsilon}_{\text{first order}}
  +\underbrace{\frac{f''(0)}{2!}\varepsilon^2}_{\text{second order}}.
  \label{eq:independence-taylor-template}
\end{equation}

The first-order term is
\begin{equation}
  f'(0)
  =\sum_{i,j}\{\pmi_P(i,j)-1\}h_{ij}=0.
  \label{eq:independence-first-term}
\end{equation}
Both parts vanish: $\pmi_P(i,j)=0$ at independence, and the probability
constraint gives $\sum_{i,j}h_{ij}=0$. Thus, positive and negative cell errors
cannot create a linear change in MI around the independence table.

To state the surviving term, define the row and column changes
$h_{i+}=\sum_jh_{ij}$ and $h_{+j}=\sum_ih_{ij}$. Differentiating the entropy
form of MI twice gives
\begin{equation}
  f''(0)
  =\sum_{i,j}\frac{h_{ij}^2}{p_{ij}}
   -\sum_i\frac{h_{i+}^2}{p_{i+}}
   -\sum_j\frac{h_{+j}^2}{p_{+j}}.
  \label{eq:independence-second-derivative}
\end{equation}
Because $p_{ij}=p_{i+}p_{+j}$, define the part of each cell change that cannot
be explained by changed margins as
\begin{equation}
  a_{ij}=h_{ij}-p_{+j}h_{i+}-p_{i+}h_{+j}.
  \label{eq:independence-association-residual}
\end{equation}
Equation~\eqref{eq:independence-second-derivative} can then be written as
\begin{equation}
  f''(0)=\sum_{i,j}\frac{a_{ij}^2}{p_{i+}p_{+j}}.
  \label{eq:independence-quadratic-form}
\end{equation}
Substituting the zeroth-, first-, and second-order terms into
Equation~\eqref{eq:independence-taylor-template} gives
\begin{equation}
  \boxed{
  I(P_\varepsilon)
  =\underbrace{0}_{\text{zeroth and first order}}
  +\underbrace{\frac{\varepsilon^2}{2}
  \sum_{i,j}\frac{a_{ij}^2}{p_{i+}p_{+j}}
  }_{\text{second order}}
  +o(\varepsilon^2)
  }.
  \label{eq:appendix-mi-second-order}
\end{equation}
The remainder $o(\varepsilon^2)$ is negligible relative to
$\varepsilon^2$ as $\varepsilon$ approaches zero.

For a multinomial sample, take $\varepsilon h=\widehat P-P$. An empirical cell
proportion has standard deviation $\sqrt{p_{ij}(1-p_{ij})/n}$, so each
$\varepsilon h_{ij}$ is typically proportional to $n^{-1/2}$. The corresponding
row- and column-adjusted errors $\varepsilon a_{ij}$ have the same scale.
Equation~\eqref{eq:appendix-mi-second-order} squares these errors, making
$\Ihat(P)$ proportional to $n^{-1}$ under independence. Multiplying by $2n$
therefore produces a statistic with a non-vanishing scale:
\begin{equation}
  2n\Ihat(P)
  \xrightarrow{d}
  \chi^2_{(r-1)(c-1)}.
\end{equation}
This is the usual likelihood-ratio, or $G$, test of independence under a fixed
alphabet and the standard positive-support large-sample conditions
\citep{wilks1938,brillinger2004}. The chi-squared reference describes the squared cell
errors in Equation~\eqref{eq:appendix-mi-second-order}. Expanded Welch instead
standardises a signed first-order MI error on the $n^{-1/2}$ scale. At
independence that signed term is absent, so its Student reference is not the
chi-squared reference in another form.

The same likelihood-ratio connection appears for conditional information
measures. For finite Markov chains, \citet{barnett2012} relate transfer entropy
to a log-likelihood ratio with an asymptotic chi-squared reference; the wider
information-theoretic setting is reviewed by \citet{bossomaier2016}. This is an
adjacent application of the same second-order likelihood principle, not a
two-sample equal-MI Welch test.

\section{What alternative constructions produce}

Using the observed margins to set
$\widehat q_{ij}=\widehat p_{i+}\widehat p_{+j}$ gives
\begin{equation}
  2nD_{\mathrm{KL}}(\widehat P\Vert\widehat Q)
  =2n\Ihat(P),
\end{equation}
which is the likelihood-ratio statistic above.
Because $\widehat Q$ is constructed from the same sample as $\widehat P$, it
is not an independent second table and cannot be inserted into the two-sample
Welch variance formula as though it supplied new information.

Repeatedly shuffling one variable preserves its observations while breaking
the pairing with the other variable. The resulting empirical distribution is
a permutation reference. Repeatedly drawing tables from
$\widehat P_X\widehat P_Y$ is a parametric bootstrap reference. A single
simulated $Q$ table adds Monte Carlo variation but does not by itself define
a calibrated null distribution.

The independent reference $Q=P_XP_Y$ is therefore mathematically valid, but
it does not create a regular two-sample MI comparison. Analytic improvement of
the sparse independence test would require a second-order finite-sample
method, rather than an immediate extension of the Expanded Welch calculation.
